\documentclass[10pt,a4paper]{article}
\usepackage[T1]{fontenc}
\usepackage[utf8]{inputenc}
\usepackage{tgtermes,tgheros,microtype}
\usepackage[a4paper,margin=18mm,headheight=19pt]{geometry}
\usepackage{xcolor,booktabs,tabularx,array,enumitem}
\usepackage[most]{tcolorbox}
\usepackage{fancyhdr,lastpage,titlesec}
\usepackage[hidelinks]{hyperref}
\usepackage{xurl,amsmath}

\definecolor{Navy}{HTML}{173753}
\definecolor{Teal}{HTML}{008F8C}
\definecolor{PaleTeal}{HTML}{E5F4F3}
\definecolor{Orange}{HTML}{F08A24}
\definecolor{PaleOrange}{HTML}{FFF1E2}
\definecolor{Red}{HTML}{BB3E35}
\definecolor{PaleRed}{HTML}{F9E9E7}
\definecolor{Green}{HTML}{2E8B57}
\definecolor{PaleGreen}{HTML}{EAF5EE}
\definecolor{Slate}{HTML}{5B7083}
\definecolor{Light}{HTML}{F3F6F8}
\definecolor{Line}{HTML}{D7E0E7}

\newcommand{\sans}{\fontfamily{qhv}\selectfont}
\newcolumntype{Y}{>{\raggedright\arraybackslash}X}
\newcolumntype{R}{>{\raggedleft\arraybackslash}X}

\hypersetup{
  pdftitle={HeiQ XMembrane Final Project Closeout -- Version 10},
  pdfauthor={HeiQ XMembrane project team},
  pdfsubject={Authoritative model-phase closeout and claim boundary},
  pdfkeywords={XMembrane, closeout, EO, ERV, capacitive electrode, claim boundary}
}

\pagestyle{fancy}
\fancyhf{}
\fancyhead[L]{\sans\footnotesize\color{Slate}HEIQ MATERIALS AG · XMEMBRANE}
\fancyhead[R]{\sans\footnotesize\color{Slate}FINAL PROJECT CLOSEOUT · V10}
\fancyfoot[L]{\sans\scriptsize\color{Slate}Model phase complete · technology not validated}
\fancyfoot[R]{\sans\scriptsize\color{Slate}Page \thepage\ of \pageref*{LastPage}}
\renewcommand{\headrulewidth}{0.45pt}
\renewcommand{\headrule}{\hbox to\headwidth{\color{Navy}\leaders\hrule height \headrulewidth\hfill}}
\renewcommand{\footrulewidth}{0pt}

\titleformat{\section}{\sans\bfseries\color{Navy}\fontsize{18}{21}\selectfont}{\thesection}{0.65em}{}
\titleformat{\subsection}{\sans\bfseries\color{Navy}\fontsize{12.5}{15}\selectfont}{\thesubsection}{0.55em}{}
\titlespacing*{\section}{0pt}{6mm}{2.5mm}
\titlespacing*{\subsection}{0pt}{4mm}{1.5mm}
\setlength{\parindent}{0pt}
\setlength{\parskip}{1.55mm}
\setlist[itemize]{leftmargin=5mm,itemsep=0.8mm,topsep=1mm}
\setlist[enumerate]{leftmargin=6mm,itemsep=1.2mm,topsep=1mm}
\renewcommand{\arraystretch}{1.18}

\newtcolorbox{mainbox}{enhanced,breakable,colback=PaleTeal,colframe=Teal,boxrule=0.7pt,arc=2mm,left=4mm,right=4mm,top=3mm,bottom=3mm}
\newtcolorbox{warningbox}{enhanced,breakable,colback=PaleOrange,colframe=Orange,boxrule=0.7pt,arc=2mm,left=4mm,right=4mm,top=3mm,bottom=3mm}
\newtcolorbox{stopbox}{enhanced,breakable,colback=PaleRed,colframe=Red,boxrule=0.7pt,arc=2mm,left=4mm,right=4mm,top=3mm,bottom=3mm}
\newtcolorbox{gobox}{enhanced,breakable,colback=PaleGreen,colframe=Green,boxrule=0.7pt,arc=2mm,left=4mm,right=4mm,top=3mm,bottom=3mm}
\newtcolorbox{plainbox}{enhanced,breakable,colback=Light,colframe=Line,boxrule=0.5pt,arc=2mm,left=4mm,right=4mm,top=3mm,bottom=3mm}
\newcommand{\eyebrow}[1]{{\sans\bfseries\color{Teal}\MakeUppercase{#1}}\par\vspace{1.2mm}}

\begin{document}
\thispagestyle{empty}
\vspace*{8mm}
\begin{tcolorbox}[enhanced,colback=Navy,colframe=Navy,boxrule=0pt,arc=0mm,left=12mm,right=12mm,top=12mm,bottom=12mm]
{\sans\bfseries\large\color{white}HEIQ XMEMBRANE PROJECT TEAM}\par
\vspace{9mm}
{\sans\bfseries\fontsize{28}{32}\selectfont\color{white}Final project closeout}\par
\vspace{2mm}
{\sans\fontsize{18}{21}\selectfont\color{white}Authoritative model-phase decision\\and claim boundary}\par
\vspace{8mm}
{\sans\large\color{white}2 September 2026\quad · \quad Version 10}
\end{tcolorbox}

\vspace{9mm}
\begin{stopbox}
{\sans\bfseries\Large\color{Red}MODEL PHASE COMPLETE}\par
{\sans\bfseries\Large\color{Red}TECHNOLOGY NOT VALIDATED}\par
\vspace{2mm}
No existing membrane, continuous electrode mechanism, module-scale performance
or net system saving has been validated. Positive energy results remain
conditional model sensitivities and are not product or commercial claims.
\end{stopbox}

\vspace{5mm}
\begin{tabularx}{\textwidth}{@{}p{43mm}Y@{}}
\textbf{Decision authority} & Current project-closeout basis; V10 governs current decisions \\
\textbf{Supersession} & Supersedes V9 for current decisions; V9 remains unchanged as historical evidence \\
\textbf{Selected model architecture} & Direct EO transport inside the ERV membrane/core \\
\textbf{Post-ERV O30} & Coupon-test boundary/control only unless explicitly reselected later \\
\textbf{Continuation} & Separate measured-data experimental and engineering phase only \\
\textbf{Release contents} & Markdown, LaTeX, compiled PDF and SHA-256 source manifest \\
\end{tabularx}

\vfill
{\sans\scriptsize\color{Slate}Controlled closeout report. Version 9 and
earlier records are preserved as historical provenance and are not modified.}

\clearpage
\section*{Executive decision}
\addcontentsline{toc}{section}{Executive decision}

This report supersedes Version~9 for all current project decisions. Version~9
documented conditional steady-state energy windows. The later
capacitive-electrode transient study materially changes their interpretation:
those windows cannot be applied to long-duration operation with purely
blocking electrodes. Where the reports differ, Version~10 governs.

At the original stored O30 terminal limit of 0.025~V, the primary replay
retains only \textbf{16.2\%} of counterfactual steady EO after one hour and
\textbf{0.68\%} over 24 hours. The former 98.5\%/4.6\% result used a separate
1.8-V current-commanded ceiling and is retained only as an optimistic upper
bound.

The modelling is sufficiently mature to stop broad hypothetical parameter
sweeps. It has delivered internally audited transport models, explicit
architecture and comparison rules, useful sensitivity coordinates and a
falsifiable coupon-test plan. It has not demonstrated a real membrane,
sustainable electrode mechanism, manufacturable module or net system saving.

\begin{mainbox}
\textbf{Final decision.} The XMembrane model phase is complete. No further
modelling is required to close this phase. Any continuation is a separately
gated experimental and engineering phase driven by measured data.
\end{mainbox}

\section*{Governing architecture}
\addcontentsline{toc}{section}{Governing architecture}

The selected modelling architecture is \textbf{direct EO transport inside the
ERV membrane/core}. Local through-plane current density and water flux are
evaluated along the supply and receiver channels; downstream conditioning
removes any remaining load.

The historical \textbf{post-ERV O30} configuration is not a selected product
or system architecture. It may be used only as a prescribed coupon-test
boundary, control or benchmark, and must be labelled that way, unless a future
controlled decision document explicitly selects another architecture. O30 is
a humidity target coordinate, not evidence that a downstream EO module is
preferred.

\begin{warningbox}
\textbf{Interpretation rule.} Historical post-ERV calculations remain
provenance. They do not override the direct EO-in-ERV architecture selected for
current modelling decisions.
\end{warningbox}

\clearpage
\section{Decisive new evidence}
\subsection{Purely blocking capacitive electrodes are transient}

For two purely blocking capacitive electrodes, charge accumulates until the
terminal-voltage limit is reached; current and incremental EO transport then
decay toward zero. The transient solver is current-commanded with an absolute
terminal-voltage clamp, unlike the voltage-controlled steady axial model. The
primary replay therefore preserves each stored case's original terminal limit:

\begin{center}\small
\begin{tabularx}{\textwidth}{@{}Yrrrrr@{}}
\toprule
\textbf{Stored case} & \textbf{Limit} & \textbf{EO 1 h} & \textbf{EO 24 h} & \textbf{Eq. time} & \textbf{Max pair V} \\
\midrule
O30 & 0.025 V & \textbf{16.2\%} & \textbf{0.68\%} & \textbf{0.17 h} & \textbf{0.478 V} \\
Full humidity duty & 0.250 V & \textbf{16.5\%} & \textbf{0.69\%} & \textbf{0.17 h} & \textbf{0.250 V} \\
\bottomrule
\end{tabularx}
\end{center}

The separate 1.8-V current-commanded replay is an optimistic upper-bound
sensitivity, not a same-voltage replay:

\begin{center}\small
\begin{tabularx}{\textwidth}{@{}Yrrrrr@{}}
\toprule
\textbf{Stored case} & \textbf{Ceiling} & \textbf{EO 1 h} & \textbf{EO 24 h} & \textbf{Eq. time} & \textbf{Max pair V} \\
\midrule
O30 upper bound & 1.800 V & \textbf{98.5\%} & \textbf{4.6\%} & \textbf{1.10 h} & \textbf{2.253 V} \\
Full-duty upper bound & 1.800 V & \textbf{92.3\%} & \textbf{5.0\%} & \textbf{1.20 h} & \textbf{1.800 V} \\
\bottomrule
\end{tabularx}
\end{center}

The percentages are relative to counterfactual continuous-current EO
displacement. The replay freezes the bulk states, gas films and installed area
while current collapses. It is a duration diagnostic, not a dynamic,
mass-balanced ERV prediction and not a revised HVAC-saving result. The
reference 10,000-F/m\textsuperscript{2} capacitance per electrode was an
unmeasured sensitivity assumption.

\begin{warningbox}
\textbf{Electrode-window limit.} The terminal-voltage clamp is not an
individual-electrode or electrode-pair electrochemical-window constraint. The
model resolves neither individual electrode potentials nor a pair stability
window. A reported pair-capacitor state may therefore exceed the terminal
clamp when membrane polarisation is negative, as in the 2.253-V O30 upper
bound. These are model states, not validated safe operating limits.
\end{warningbox}

Symmetric, charge-balanced cycling gives \textbf{zero net EO water transport
over a complete cycle} in the present constant-drag model, to numerical
precision. Water moved during one half-cycle is returned during the other.
A sustainable result requires a measured asymmetry or another physical
mechanism; it cannot be inferred from capacitance alone.

\begin{stopbox}
\textbf{Decision consequence.} Earlier steady energy windows must not be
reported as sustained performance for purely blocking electrodes.
\end{stopbox}

\subsection{The recent combined material case is synthetic}

The recent joint \(t_w=50\), \(D=1.0\times10^{-10}\)~m\textsuperscript{2}/s
run is retained as a sensitivity result, not a candidate membrane.

\begin{center}
\begin{tabularx}{0.86\textwidth}{@{}YR@{}}
\toprule
\textbf{Item} & \textbf{Result} \\
\midrule
Transport-closed rows & \textbf{4 of 8} \\
Best available eight-segment diagnostic area & \textbf{43,019.5 m\textsuperscript{2}} \\
Voltage at that diagnostic & \textbf{0.10 V} \\
Full-coupled/full-system rows & \textbf{0} \\
\bottomrule
\end{tabularx}
\end{center}

The audit passed integrity and reproducibility checks only. The joint property
vector was not measured on an identifiable membrane; the material-validation
flag is false. Area is diagnostic and the conditional energy arithmetic omits
unvalidated subsystems. The run therefore does not establish a realistic
material, practical module or system-energy claim.

\subsection{Status of earlier energy windows}

Version~9's integrated-ERV O30 result and other positive energy differences
remain useful records of what the steady model predicted under its declared
assumptions. They are historical conditional evidence. They must not be
reported as sustained performance because continuous electrodes, same-sample
membrane properties, receiver/thermal coupling and practical hardware remain
unvalidated.

\clearpage
\section{Technical reference for the frozen axial model}

The steady axial calculations are \textbf{voltage controlled}: terminal
voltage is held constant for a case, while current is not. Each segment solves
its local through-plane current density \(j(x)\), coupled to gas films and net
water flux, from
\[
V_{\mathrm{app}}=j(x)\!\left(\mathrm{ASR}_{\mathrm{mem}}(x)+
\mathrm{ASR}_{\mathrm{series}}\right)
+t_w\frac{RT}{F}\ln\!\left(\frac{a_{w,r}^{s}}{a_{w,s}^{s}}\right)
+V_{\mathrm{fixed}}.
\]
Activities, resistance and flux vary along the channel, so \(j(x)\) varies
with axial position. This is through-plane current sampled along \(x\), not a
solved in-plane current-collector current. It has equal physical magnitude on
both membrane faces; its sign reverses only when written using each face's own
outward normal.

Both channel gaps are 2.0~mm, so \(D_h=4.0\)~mm. The model has no literal
gas-film slab. At film scale factor 1, \(Sh=7.54\) gives the \textbf{derived
effective film thickness}
\[
\delta_{\mathrm{eff}}=\frac{D}{k}=\frac{D_h}{Sh}=0.5305\ \mathrm{mm}
\]
on each side. At 32/20~\textdegree C, supply/receiver
\(k\approx0.0511/0.0475\)~m/s. The 0.5305-mm value is a correlation-equivalent
boundary layer, not a physical GDL or manufactured film.

The inverse-design anchor is hypothetical: \(t_w=100.0608\),
\(D=1.5979\times10^{-14}\)~m\textsuperscript{2}/s,
\(\sigma=23.1719\)~S/m and dry \(L=36.9137\)~\(\mu\)m. It must not be called
a realistic membrane. The \(t_w=50\), \(D=10^{-10}\)~m\textsuperscript{2}/s
vector is likewise synthetic.

\clearpage
\section{Final claim boundary}
\eyebrow{What may and may not be stated}

The following statements are supported:
\begin{itemize}
  \item the model phase is complete;
  \item direct EO inside the ERV is the selected modelling architecture;
  \item the models identify conditional sensitivities and experimental hypotheses;
  \item purely blocking capacitive electrodes provide finite startup
  displacement, not demonstrated continuous EO duty;
  \item the original-voltage O30 replay retains 16.2\% after one hour and
  0.68\% over 24 hours; the 1.8-V replay is only an optimistic
  current-commanded upper bound;
  \item the \(t_w=50\), \(D=10^{-10}\)~m\textsuperscript{2}/s run is a
  synthetic, partially closed sensitivity screen.
\end{itemize}

The following statements are \textbf{not supported} and must not be made:
\begin{itemize}
  \item an existing or procurable membrane meets the model property bundle;
  \item continuous EO transport has been demonstrated;
  \item a Faradaic, pseudocapacitive or rectifying electrode architecture is validated;
  \item membrane, module, receiver, current collector or full-system
  performance has been experimentally validated;
  \item a net HVAC energy saving, commercial advantage, lifetime, cost or
  manufacturability target has been achieved.
\end{itemize}

\begin{stopbox}
\textbf{Controlled closeout statement.} The model phase is complete. No
existing membrane, continuous electrode mechanism, module-scale performance or
net system saving has been validated. All positive energy results remain
conditional model sensitivities and must not be presented as product or
commercial claims.
\end{stopbox}

\section{Accepted successor phase}
\eyebrow{Measured evidence before another design claim}

Only the following work is accepted if the technology is continued:

\begin{enumerate}
  \item \textbf{Demonstrate a sustainable electrode mechanism.} Measure a
  Faradaic, pseudocapacitive, rectifying, flow-switched or other charge/water-
  asymmetric mechanism. Quantify separate electrode potentials, reaction
  products, kinetics, overpotential, energy recovery, self-discharge and lifetime.
  \item \textbf{Complete MAT-01--MAT-09 on the same identifiable sample or
  batch.} Measure sorption/hysteresis, mobility/permeance, EO drag, wet
  through-plane conductivity, membrane/contact ASR, thickness/swelling/
  mechanics, coupled net water flux, interface compatibility and provenance.
  \item \textbf{Run a controlled-flow coupon test.} Use matched EO-off/EO-on,
  ABBA polarity reversal, two-sided water balance, pressure control, film/gap
  repeats and complete charge/discharge cycles. Post-ERV O30 may be used only
  as the declared boundary/control.
  \item \textbf{Close the dynamic and full-system ledgers.} Couple transient
  membrane inventory and gas streams, receiver operation, heat, pressure drop,
  fans, power electronics, heat rejection, leakage and cross-contamination.
  \item \textbf{Apply product constraints.} Set accepted limits for active
  area, volume, pressure drop, current collection, cost, manufacturability,
  safety, fouling, lifetime and end-of-life handling.
  \item \textbf{Replay once with measured inputs and hold a new gate review.}
  Do not restart broad hypothetical sweeps. Any new model or product claim
  requires a controlled successor report.
\end{enumerate}

\clearpage
\section{Closeout register}

\begin{center}
\begin{tabularx}{\textwidth}{@{}Yp{57mm}@{}}
\toprule
\textbf{Workstream} & \textbf{Closeout status} \\
\midrule
Mathematical/model development & \textcolor{Green}{\textbf{Complete for this phase}} \\
Sensitivity and architecture screening & \textcolor{Green}{\textbf{Complete for this phase}} \\
Real material qualification & \textcolor{Orange}{\textbf{Not validated; successor phase}} \\
Sustainable electrode operation & \textcolor{Orange}{\textbf{Not validated; successor phase}} \\
Coupled module/full-system closure & \textcolor{Orange}{\textbf{Not validated; successor phase}} \\
Product/commercial decision & \textcolor{Red}{\textbf{No claim permitted}} \\
\bottomrule
\end{tabularx}
\end{center}

\begin{gobox}
\textbf{Administrative closeout.} The scientific modelling phase can be
closed without executing the successor experimental programme. The open items
are accepted future-phase gates, not incomplete deliverables of this model
phase.
\end{gobox}

\subsection*{Release verification}

The regenerated capacitive audit passes \textbf{16 of 16} checks and its
targeted module passes \textbf{9 of 9} tests. The corrected targeted release
suite passes \textbf{63 tests}; the full suite passes \textbf{313 tests with 1
skipped}. The sole skip is the optional design-frontier smoke test because its
output is absent. These are code and integrity results, not technology
validation; the governing package result remains the V10 status file.

\path{scripts/run_final_delivery_check_v1.py} is a superseded historical
checker. It predates the September evidence and fails on stale August report
and nested-manifest assumptions, so its result is not a V10 acceptance signal.
The sole governing release status is the new fail-closed package verification
at
\path{deliverables/current/2026-09-02_v10/verification/V10_STATUS.json},
produced by:

\begin{plainbox}
{\small\ttfamily python3 scripts/package\_final\_closeout\_v10.py\par
\hspace*{5mm}--expected-source-revision ``\$(git rev-parse HEAD)''}
\end{plainbox}

This report does not claim that package verification passes before that status
file has been generated successfully.

\section{Controlled evidence basis}

\begin{itemize}
  \item \path{reports/final_team_decision_report_2026-08-31/HeiQ_XMembrane_Final_Result_and_Decision_Report_2026-08-31_v9.tex} --- preserved historical V9 report;
  \item \path{docs/CAPACITIVE_ELECTRODE_TRANSIENT_STUDY_2026-09-01.md} and
  \path{outputs/capacitive_electrode_study_v1/} --- transient study and audit;
  \item \path{outputs/direct_active_erv_tw50_d1e10_v1/study/} --- combined
  synthetic sensitivity and audit;
  \item \path{outputs/integrated_evidence_v1/material_data_requirements.csv}
  --- MAT-01--MAT-09 qualification schema;
  \item \path{src/heiq/wp1/axial/local_state.py},
  \path{src/heiq/wp1/airside/properties.py} and
  \path{heiq_pre_receiver_material_analyses/heiq_pre_receiver_model.py} ---
  local voltage/current and gas-film model basis;
  \item \path{reports/final_project_closeout_2026-09-02/source_manifest.json}
  --- SHA-256 bindings for this release and decision-critical sources.
\end{itemize}

\vfill
{\sans\scriptsize\color{Slate}End of controlled Version~10 closeout. V9 and
earlier files remain available as historical provenance. V10 governs current
project decisions.}
\end{document}
